\section{Detector-response-mismatch closure}
\label{app:response-mismatch}

The truth-shape and hidden-variable closures in
\S\ref{sec:validation} test whether the estimator recovers deliberately
modified event populations. They do not directly test a different question:
whether the inferred truth distribution is biased when the detector response
used to generate the pseudo-data disagrees with the nominal response supplied
to OmniFold. Kevin McFarland suggested this distinction and proposed stressing
the hadronic record, especially for quasielastic events, to inspect the
resulting bias at low $W$ (private communication).

\paragraph{Record-only test implemented in the pipeline.}
The first-stage diagnostic begins with the ordinary five-dimensional closure
sample. After the pseudo-data arrays have been copied, and before background
mixing or classifier training, it replaces only the reconstructed pseudo-data
record
\begin{equation}
  \Eavail^{\mathrm{pseudo}} \longrightarrow
  (1+\epsilon)\Eavail^{\mathrm{pseudo}} .
\end{equation}
The nominal simulated reconstructed sample, all truth coordinates, event
weights, and event membership remain fixed. Consequently the pseudo-data are
generated with a deliberately incompatible conditional response while the
unfold still sees the nominal response. Both signs of $\epsilon$ must be run
at an amplitude fixed before inspecting the result. The amplitude should be
tied to a documented hadronic-response variation or to a quantified assumption;
the inherited factor applied in the reconstructed-$\Eavail$ definition does not
by itself supply that justification.

The diagnostic writes the full five-dimensional unfolded/reference residual,
the one-dimensional $\Eavail$ and $W$ projections, and an
$(\Eavail,W)$ ratio map. The primary observable is the signed bias resolved in
$W$, with the $W<\SI{1.1}{GeV}$ region reported separately from the high-$W$
tail. The global integral is a control rather than the decision statistic.
The implementation refuses real-data mode, truth reweighting, bootstrapping,
or a non-five-dimensional axis list, and marks every output
\texttt{NONQUOTABLE-DIAGNOSTIC}.

\paragraph{Physically correlated follow-up.}
The record-only arm is intentionally cheap and adversarial, not a detector
model. A more faithful follow-up would select quasielastic events, shift the
leading-nucleon energy up and down, and then recompute reconstructed $\Eavail$,
$q_{3}$, $W$, and any affected selection from the altered low-level event
record. The current scalar event-loop output does not preserve enough
information to perform that recomputation after the fact. This arm therefore
requires a dedicated diagnostic event-loop output and is deferred unless the
record-only test shows a material bias or the collaboration requests the
stronger model.

\paragraph{Status and interpretation.}
This appendix records a validation protocol, not a result or an uncertainty
component. No response-mismatch artifact has been produced, no pass threshold
has been adopted, and the test does not alter the frozen central value or the
current covariance construction. A small observed bias would support the
stability of the $(\Eavail,W)$ interpretation only over the tested response
envelope; a large bias would trigger a model/uncertainty decision rather than
being silently absorbed as a non-closure error.
